\documentclass[12pt,a4paper]{article}
\usepackage[utf8]{inputenc}
\usepackage[spanish]{babel}
\usepackage{amsmath,amssymb,amsfonts}
\usepackage{geometry}
\usepackage{hyperref}
\geometry{margin=2.5cm}

\title{\textbf{Constantes Matemáticas en Geometría Analítica}}
\author{Compendio de Constantes Fundamentales}
\date{\today}

\begin{document}

\maketitle

\section*{Introducción}
La geometría analítica conecta el álgebra y la geometría euclidiana mediante el uso de sistemas de coordenadas cartesianas. Las constantes en geometría analítica definen invariantes de cónicas (excentricidades), razones trigonométricas notables y pendientes de pendientes geométricas clave.

\section{Excentricidad de la Parábola ($e_{\text{parábola}}$)}

\subsection*{1. Significado, Símbolo y Explicación}
La excentricidad $e$ es una medida del grado de desviación de una cónica respecto a una circunferencia perfecta. Para cualquier parábola, por definición geométrica del lugar geométrico (distancia constante al foco y a la directriz), la excentricidad es exactamente igual a la unidad:
\begin{equation*}
e_{\text{parábola}} = \frac{d(P, F)}{d(P, D)} = 1
\end{equation*}
Determina la forma geométrica en reflectores parabólicos, antenas satelitales y trayectorias de escape orbitales de energía nula.

\subsection*{2. Valor Típico en Ingeniería y Ciencia}
\begin{equation*}
e_{\text{parábola}} = 1.0
\end{equation*}

\subsection*{3. Valor Exacto a 15 Decimales}
\begin{equation*}
e_{\text{parábola}} = 1.000000000000000
\end{equation*}

\subsection*{4. Valor Exacto a 100 Decimales}
\begin{equation*}
e_{\text{parábola}} = 1.0000000000000000000000000000000000000000000000000000000000000000000000000000000000000000000000000000
\end{equation*}
\textbf{Margen de error:} Exacto por definición geométrica.

\vspace{0.8cm}
\hrule
\vspace{0.8cm}

\section{Pendiente de la Bisectriz del Primer Cuadrante ($m_{1}$)}

\subsection*{1. Significado, Símbolo y Explicación}
En el plano cartesiano $\mathbb{R}^2$, la bisectriz del primer y tercer cuadrante es la recta de ecuación $y = x$. Su pendiente, denotada por $m_1$, es igual a la tangente del ángulo de inclinación de $45^\circ$ ($\pi/4$ rad):
\begin{equation*}
m_1 = \tan(45^\circ) = \tan\left(\frac{\pi}{4}\right) = 1
\end{equation*}
Representa la proporcionalidad directa $1:1$ entre los ejes horizontal y vertical.

\subsection*{2. Valor Típico en Ingeniería y Ciencia}
\begin{equation*}
m_1 = 1.0
\end{equation*}

\subsection*{3. Valor Exacto a 15 Decimales}
\begin{equation*}
m_1 = 1.000000000000000
\end{equation*}

\subsection*{4. Valor Exacto a 100 Decimales}
\begin{equation*}
m_1 = 1.0000000000000000000000000000000000000000000000000000000000000000000000000000000000000000000000000000
\end{equation*}
\textbf{Margen de error:} Exacto por definición analítica.

\end{document}
